% !TEX root = ../Projektdokumentation.tex
\section{Einleitung}
\label{sec:Einleitung}

\subsection{Vorstellung des DKFZ und der ODCF}
\label{sec:Vorstellung_DKFZ}
Das Deutsche Krebsforschungszentrum (\ac{DKFZ}) in Heidelberg ist die größte
biomedizinische Forschungseinrichtung Deutschlands und betreibt
Spitzenforschung im Bereich der Krebsentstehung und -bekämpfung.
Innerhalb des DKFZ stellt die \ac{ODCF} (Omics IT and Data Management Core
Facility) den Forschungsgruppen eine komplexe IT-Infrastruktur zur
Verarbeitung großer Mengen genomischer Daten bereit.
Dazu zählen unter anderem Hochleistungsrechner, Storage-Systeme sowie eine
Virtualisierungsinfrastruktur auf Basis von \ac{PVE}.

\subsection{Projektumfeld}
\label{sec:Projektumfeld}
Das Projekt wird in den Räumen der \ac{ODCF} am \ac{DKFZ} in Heidelberg
durchgeführt, wo der Prüfling seine Ausbildung zum Fachinformatiker
Systemintegration absolviert.
Auftraggeber und fachlicher Betreuer ist Frank Thommen (\ac{ODCF}).
Folgende Beteiligte und Schnittstellen sind für die Projektdurchführung
relevant:
\begin{itemize}
    \item Betreuer / Ausbilder: verantwortlich für die fachliche Begleitung und die Abnahme des Projekts
    \item \ac{ODCF}-Systemadministratoren: Betreiber der \ac{PVE}-Umgebung und Übergabeempfänger des fertigen \ac{PBS}
    \item Forschungsgruppen des \ac{DKFZ}: Nutzer der virtualisierten Dienste (indirekt betroffen)
    \item Netzwerkteam: zuständig für \ac{VLAN}-Konfiguration und Firewall-Freischaltungen
\end{itemize}

\subsubsection{Ist-Situation}
Die \ac{ODCF} betreibt eine \ac{PVE}-Cluster-Umgebung, die eine Vielzahl
\acp{VM} und \acp{LXC} für Forschungs- und Infrastrukturaufgaben
bereitstellt.
Eine grundlegende Backup-Infrastruktur ist vorhanden: Sicherungen der
virtuellen Maschinen werden aktuell auf Hypervisor-Ebene als Snapshots
erstellt und anschließend auf Tape (Bandlaufwerk) gesichert.
Ein zentralisierter, dedizierter \acl{PBS} für die Verwaltung, Überwachung
und effiziente inkrementelle Datensicherung aller \acp{VM} und Container
existiert jedoch nicht.
Dies führt zu eingeschränkter Wiederherstellbarkeit, fehlender zentraler
Verwaltung der Backup-Jobs sowie begrenzten Möglichkeiten zur Überprüfung
der Backup-Integrität.

\subsection{Projektziel}
\label{sec:Projektziel}
Ziel des Projekts ist die Einrichtung eines dedizierten \acl{PBS} sowie
dessen vollständige Integration in die bestehende \ac{PVE}-Umgebung der
\ac{ODCF}.
Der \ac{PBS} ergänzt die vorhandene Tape-basierte Sicherung um eine
zentrale, effiziente und inkrementelle Backup-Lösung auf Disk-Ebene.
Nach Projektabschluss sollen alle relevanten \acp{VM} und Container
automatisiert über den \ac{PBS} gesichert werden.
Konkret sind folgende Ziele zu erreichen:
\begin{itemize}
    \item Installation und Konfiguration eines \acl{PBS} auf dedizierter Hardware
    \item Anbindung des \ac{PBS} an den bestehenden \ac{PVE}-Cluster der \ac{ODCF}
    \item Einrichtung von Backup-Jobs mit definierten Zeitplänen und Aufbewahrungsrichtlinien (Retention Policies: tägliche, wöchentliche, monatliche Sicherungen)
    \item Konfiguration von Benachrichtigungen und Monitoring für Backup-Jobs
    \item Erstellung einer Betriebsdokumentation und Übergabe an das \ac{ODCF}-Team
\end{itemize}

\subsection{Projektbegründung}
\label{sec:Projektbegruendung}
Die bisherige, ausschließlich auf Tape basierende Sicherung bietet zwar
einen Schutz vor Datenverlust, ist aber für die operativen Anforderungen
einer Forschungsinfrastruktur nicht ausreichend.
Restores von einzelnen \acp{VM} oder gar einzelnen Dateien sind aufwändig,
inkrementelle Sicherungen sind nicht möglich und eine zentrale Überwachung
der Backup-Jobs fehlt.
Mit einem dedizierten \acl{PBS} können Sicherungen deutlich häufiger,
effizienter (Deduplizierung, inkrementelle Backups) und mit deutlich
kürzeren Wiederherstellungszeiten durchgeführt werden.
Gleichzeitig erlaubt der \ac{PBS} eine zentrale Verwaltung,
Integritätsprüfung und Benachrichtigung über fehlgeschlagene Jobs.

\subsection{Projektabgrenzung}
\label{sec:Projektabgrenzung}
Im Rahmen des Abschlussprojekts werden ausschließlich der \acl{PBS} und
seine Anbindung an die bestehende \ac{PVE}-Umgebung der \ac{ODCF}
betrachtet.
Die bestehende Tape-basierte Sicherung bleibt parallel in Betrieb und
wird durch das Projekt nicht abgelöst.
Themen wie Off-Site-Replikation des \ac{PBS}-Datastores, Desaster-Recovery-
Konzepte über das Rechenzentrum hinaus oder die produktive Migration aller
Forschungsdienste sind nicht Teil dieses Projekts.
Ebenfalls außerhalb des Umfangs liegt die Anpassung der einzelnen
Anwendungen innerhalb der \acp{VM}.

\subsection{Projektschnittstellen}
\label{sec:Projektschnittstellen}
Das System interagiert mit folgenden Komponenten der \ac{ODCF}-Infrastruktur:
\begin{itemize}
    \item \acl{PVE}-Cluster: Quelle der zu sichernden \acp{VM} und Container
    \item Netzwerk-Infrastruktur (\ac{VLAN}, Firewall): Anbindung des \ac{PBS}
    \item E-Mail-Server / Monitoring: Versand von Benachrichtigungen über Backup-Jobs
    \item Bestehendes Tape-Backup: bleibt parallel als zweite Sicherungsebene erhalten
\end{itemize}
Die direkten Nutzer des \ac{PBS} sind die \ac{ODCF}-Systemadministratoren,
die nach der Übergabe für den laufenden Betrieb und die regelmäßige
Überprüfung der Sicherungen zuständig sind.

\subsection{Voraussetzungen für das Verständnis}
\label{sec:Voraussetzungen}
Für das Verständnis der folgenden Kapitel werden einige Technologien kurz
vorgestellt.

\subsubsection{\acl{PVE}}
\acl{PVE} ist eine quelloffene Virtualisierungsplattform auf Basis von
Debian Linux, die sowohl \acp{VM} (KVM) als auch \acp{LXC} verwalten kann.
\ac{PVE} wird in der \ac{ODCF} als Cluster betrieben, sodass mehrere
Hypervisor-Knoten zentral verwaltet werden.

\subsubsection{\acl{PBS}}
Der \ac{PBS} ist die hauseigene Backup-Lösung von Proxmox.
Er ist eng mit \ac{PVE} integriert und unterstützt inkrementelle
Sicherungen, clientseitige Deduplizierung sowie Integritätsprüfungen
über Prüfsummen.
Der \ac{PBS} kann mehrere Datastores verwalten, auf denen die Sicherungen
abgelegt werden.

\subsubsection{\acl{ZFS}}
\acl{ZFS} ist ein modernes Dateisystem mit integrierter
Volume-Verwaltung. Es bietet Prüfsummen für alle Daten, Snapshots,
Kompression und einfache Erweiterbarkeit. Für den Datastore des \ac{PBS}
ist \ac{ZFS} aufgrund seiner Datenintegritätsmerkmale gut geeignet.

\subsubsection{Retention Policies}
Retention Policies legen fest, wie viele und welche Sicherungen über
welchen Zeitraum vorgehalten werden. Üblich sind gestaffelte
Aufbewahrungen (z.\,B. tägliche Backups der letzten Woche, wöchentliche
des letzten Monats, monatliche des letzten Jahres). Der \ac{PBS} setzt
diese Vorgaben über sogenannte \emph{Prune}-Regeln automatisch um.
